% !TEX root = ../algorithms.tex

	\section{Ранжирование списка во внешней памяти. Линейное программирование}
	
	\subsection{Задача}
	Дан список, хранящийся в массиве $A[0..N-1]$ в произвольном порядке.
	У каждого узла $x$ есть поля
	\[
	x.\texttt{next},
	\qquad
	x.w.
	\]
	Также в реализации удобно хранить поле $x.\texttt{id}$, причём изначально
	\[
	A[i].\texttt{id}=i.
	\]
	
	\medskip
	\noindent
	\begin{Remark}{}{}
		Ниже в рассуждениях мы иногда будем писать $\operatorname{pred}(x)$ для предшественника вершины $x$ в текущем списке.
		Это \emph{не поле структуры данных}, а только удобное математическое обозначение для анализа и для псевдокода.
	\end{Remark}
	
	Исходная задача состоит в том, чтобы для каждой вершины $x$ определить её номер в списке, считая от головы:
	\[
	1,2,\dots,N.
	\]
	Эквивалентно, нужно упорядочить вершины массива $A$ в порядке их следования вдоль списка.
	
	Для решения этой задачи удобно рассмотреть более общую подзадачу:
	для каждой вершины $x$ вычислить сумму весов на префиксе списка до неё, то есть
	\[
	\mathrm{pref}(x)=\sum_{y \text{ от головы списка до } x} y.w.
	\]
	
	После работы алгоритма в поле $x.w$ будет записано значение $\mathrm{pref}(x)$.
	
	Если изначально у всех вершин $w=1$, то $\mathrm{pref}(x)$ есть просто номер вершины $x$ в списке, считая от головы:
	\[
	1,2,\dots,N.
	\]
	
	\subsection{Идея сжатия списка}
	
	Пусть $S$ --- множество удаляемых узлов, такое что никакие два соседних по списку узла не удаляются одновременно:
	\[
	\forall x\in S\quad \operatorname{pred}(x)\notin S,
	\qquad |S|=\Theta(N).
	\]
	
	Тогда все узлы из $S$ можно удалить \emph{независимо друг от друга}:
	для каждого $x\in S$ выполняем
	\[
	\operatorname{pred}(x).\texttt{next}=x.\texttt{next},
	\qquad
	\operatorname{pred}(x).w \mathrel{+}= x.w.
	\]
	
	После этого каждый оставшийся узел хранит суммарный вес некоторого подряд идущего блока исходных вершин.
	
	На получившемся более коротком списке мы рекурсивно решаем \emph{ту же самую задачу}: вычисляем префиксные суммы этих объединенных весов.
	
	После возврата из рекурсии удалённые вершины нужно восстановить и дописать для них ответы
	
	\subsection{Как выбрать множество \(S\)}
	
	Каждой вершине независимо присваиваем случайный бит
	\[
	x.r\in\{0,1\},
	\qquad
	\Pr[x.r=0]=\Pr[x.r=1]=\frac12.
	\]
	
	После этого выбираем
	\[
	x\in S
	\iff
	x.r=1 \ \text{и}\ \operatorname{pred}(x).r=0.
	\]
	
	То есть в $S$ попадают вершины, которые являются ``правым концом'' шаблона $01$ вдоль списка.
	
	\subsubsection*{Почему это корректно}
	
	Если $x\in S$, то по определению $\operatorname{pred}(x).r=0,\qquad x.r=1$.

	Значит, $\operatorname{pred}(x)$ не может одновременно лежать в $S$, потому что для этого у него должен был бы быть бит $1$.
	
	Следовательно, $\forall x\in S\colon \operatorname{pred}(x)\notin S$.
	
	\subsubsection*{Размер множества $S$}
	
	Для внутренней вершины вероятность попасть в $S$ равна
	\[
	\Pr[x\in S]=\Pr[\operatorname{pred}(x).r=0 \land x.r=1]=\frac14.
	\]
	
	Поэтому $\mathbb{E}|S| \approx \frac{N}{4}$,

	а после одного шага остаётся примерно $N-|S| \approx \frac{3N}{4}$
	вершин.
	
	\subsection{Рекурсивная схема алгоритма}
	
	Следующий псевдокод --- это \emph{не} реализация во внешней памяти, а только удобная запись логики алгоритма, где разрешено использовать обозначение $\operatorname{pred}(x)$.
	
	\begin{algorithm}[H]
		\caption{\textsc{PrefixSums-Conceptual}$(A)$}
		\begin{algorithmic}[1]
			\Require Текущий список $A$
			\Ensure Для каждой вершины $x$ значение $x.w$ равно сумме весов на префиксе до $x$
			
			\If{$|A| \le 1$}
			\State \Return $A$
			\EndIf
			
			\State Для каждой вершины $x \in A$ независимо выбрать случайный бит $x.r \in \{0,1\}$
			\State $S = \{\, x \in A \mid \operatorname{pred}(x)\ \text{существует},\ \operatorname{pred}(x).r = 0,\ x.r = 1 \,\}$
			
			\Comment{Сжимаем список, удаляя вершины из $S$}
			\ForAll{$x \in S$}
			\State $p = \operatorname{pred}(x)$
			\State $p.w = p.w + x.w$
			\State $p.\texttt{next} = x.\texttt{next}$
			\EndFor
			
			\State $A' = A \setminus S$
			\State \Call{PrefixSums-Conceptual}{$A'$}
			
			\Comment{Восстанавливаем удалённые вершины}
			\ForAll{$x \in S$ в порядке, обратном удалению}
			\State $p = \operatorname{pred}(x)$
			\State $p.\texttt{next} = x$
			\State $p.w = p.w - x.w$ \Comment{теперь $p.w$ = префиксная сумма до $p$}
			\State $x.w = x.w + p.w$ \Comment{теперь $x.w$ = префиксная сумма до $x$}
			\EndFor
			
			\State \Return $A$
		\end{algorithmic}
	\end{algorithm}
	
	\subsection{Реализация через сортировки во внешней памяти}
	
	Главная цель --- заменить случайные обращения к указателям на несколько сортировок и линейных проходов.
	
	\begin{Remark}{}{}
		Ниже для краткости не выписана отдельная обработка головы списка, хвоста списка и значения \textsc{null} в поле \texttt{next}.
		В реальной реализации эти пограничные случаи нужно обрабатывать отдельно
		(например, специальными фиктивными вершинами или отдельными ветками в коде).
	\end{Remark}
	\subsubsection*{Храним идентификаторы}
	
	У каждого узла есть поле
	\[
	x.\texttt{id},
	\qquad
	A[i].\texttt{id}=i.
	\]
	
	Это позволяет после перестановок и сортировок восстановить ``исходную'' вершину.
	
	\begin{algorithm}[H]
		\caption{\textsc{EM\_PrefixSums}$(A)$}
		\begin{algorithmic}[1]
			\Require Массив вершин $A$, отсортированный по $\texttt{id}$
			\Ensure Для каждой вершины $x$ значение $x.w$ равно сумме весов на префиксе до $x$
			
			\If{$|A| \le 1$}
			\State \Return $A$
			\EndIf
			
			\State $(\widehat A, S) = \Call{Compress}{$A$}$
			\Comment{$\widehat A$ --- сжатый список, $S$ --- удалённые вершины}
			
			\State $\widehat A = \Call{EM\_PrefixSums}{\widehat A}$
			\Comment{рекурсивно решаем ту же задачу на объединенных весах}
			
			\State $A = \Call{Restore}{\widehat A, S, |A|}$
			\State \Return $A$
		\end{algorithmic}
	\end{algorithm}
	
	\subsubsection*{Сортируем по \texttt{next}}
	
	Пусть массив $A$ отсортирован по \texttt{id}. Тогда строим
	\[
	\widetilde A = \Sort(\{x\in A\}\text{ по }x.\texttt{next}).
	\]
	
	После такой сортировки запись $\widetilde A[i]$ стоит напротив той вершины, в которую она указывает.
	Поэтому при одновременном просмотре массивов $A$ и $\widetilde A$ можно сопоставить вершину и её предшественника.
	
	\subsubsection*{Находим удаляемые вершины и сразу обновляем предшественников}
	
	При линейном проходе по $A$ и $\widetilde A$:
	
	\[
	\text{если } \widetilde A[i].r=0 \text{ и } A[i].r=1,
	\text{ то удаляем } A[i].
	\]
	
	При этом выполняются обновления:
	\[
	\widetilde A[i].w \mathrel{+}= A[i].w,
	\qquad
	\widetilde A[i].\texttt{next}=A[i].\texttt{next},
	\]
	а сама удалённая вершина сохраняется в массив $S$.
	
	\subsubsection*{Убираем ``дубликаты'' после перенаправления указателей}
	
	После удаления вершины её предшественник начинает указывать в того же преемника, в которого раньше указывала удалённая вершина.
	Поэтому после повторной сортировки по \texttt{next}
	\[
	\widetilde A = \Sort(\widetilde A\text{ по } \texttt{next})
	\]
	появляются соседние элементы с одинаковым полем \texttt{next}.
	
	Из каждой такой пары надо оставить только живую вершину, а удалённую выбросить.
	Это делается линейным проходом по отсортированному массиву:
	если у двух соседних элементов одинаковый \texttt{next}, то удаляется тот, у которого $r=1$.
	
	После этого выжившие вершины упаковываются в плотный массив длины $N-|S|$.
	
	\begin{algorithm}[H]
		\caption{\textsc{Compress}$(A)$}
		\begin{algorithmic}[1]
			\Require Массив $A[0..N-1]$, отсортированный по $\texttt{id}$
			\Ensure Сжатый массив $\widehat A$ и массив удалённых вершин $S$
			
			\State Для каждой вершины $x \in A$ независимо выбрать случайный бит $x.r \in \{0,1\}$
			\State $\widetilde A = \Sort(A \text{ по } x.\texttt{next})$
			\State $S = [\ ]$
			
			\Comment{Находим вершины из шаблона $01$ и сразу переносим их вес в предшественника}
			\For{$i=0$ \textbf{to} $N-1$}
			\If{$\widetilde A[i].r = 0$ \textbf{and} $A[i].r = 1$}
			\State $\widetilde A[i].w = \widetilde A[i].w + A[i].w$
			\State $\widetilde A[i].\texttt{next} = A[i].\texttt{next}$
			\State $S.\textsc{push\_back}(A[i])$
			\EndIf
			\EndFor
			
			\State $\widetilde A = \Sort(\widetilde A \text{ по } x.\texttt{next})$
			
			\Comment{Упаковываем только живые вершины}
			\State $B = [\ ]$
			\State $i = 0$
			\While{$i < N$}
			\If{$i+1 < N$ \textbf{and} $\widetilde A[i].\texttt{next} = \widetilde A[i+1].\texttt{next}$}
			\If{$\widetilde A[i].r = 1$}
			\State swap$(\widetilde A[i], \widetilde A[i+1])$
			\Comment{после swap живая вершина стоит первой}
			\EndIf
			\State $B.\textsc{push\_back}(\widetilde A[i])$
			\State $i = i + 2$
			\Else
			\State $B.\textsc{push\_back}(\widetilde A[i])$
			\State $i = i + 1$
			\EndIf
			\EndWhile
			
			% Изменено: явно прописали перенумерацию next после упаковки.
			\Comment{Переопределяем \texttt{next}: старые \texttt{id} заменяются на новые позиции в плотном массиве}
			\State $\widetilde B = \Sort(B \text{ по } x.\texttt{next})$
			\State $i = 0$; $j = 0$; $k = 0$
			\While{$j < N-|S|$}
			\While{$k < |S|$ \textbf{and} $i = S[k].\texttt{id}$}
			\State $i = i + 1$
			\State $k = k + 1$
			\EndWhile
			\State $\widetilde B[j].\texttt{next} = j$
			\State $i = i + 1$
			\State $j = j + 1$
			\EndWhile
			
			% Изменено: явно сказали, что id и w не переинициализируются.
			\State $\widehat A = \Sort(\widetilde B \text{ по } x.\texttt{id})$
			\Comment{$\texttt{id}$ сохраняются исходными, веса $w$ уже агрегированы}
			\State \Return $(\widehat A, S)$
		\end{algorithmic}
	\end{algorithm}
	
	\subsubsection*{Перенумерация и рекурсия}
	
	После упаковки нужно переопределить индексы \texttt{next} под новый массив,
	затем отсортировать вершины по \texttt{id} и рекурсивно решить задачу уже для массива размера $N-|S|$:
	\[
	\widehat A = \Sort(\{x\in A[0..N-|S|-1]\}\text{ по }x.\texttt{id}).
	\]
	
	\subsubsection*{Восстановление}
	
	Предполагается, что:
	\begin{itemize}
		\item $S$ хранится в порядке возрастания \texttt{id};
		\item $\widehat A$ --- результат рекурсивного решения на сжатом массиве длины $N-|S|$.
	\end{itemize}
	
	Сначала сортируем удалённые и оставшиеся вершины по полю \texttt{next}:
	\[
	\widetilde S = \Sort(\{x\in S\}\text{ по }x.\texttt{next}),
	\qquad
	\widetilde A = \Sort(\{x\in \widehat A[0..N-|S|-1]\}\text{ по }x.\texttt{next}).
	\]
	
	После этого выполняем линейный проход:
	\begin{algorithm}[H]
		\caption{\textsc{Restore}$(\widehat A, S, N)$}
		\begin{algorithmic}[1]
			\State $\widetilde S = \Sort(S \text{ по } x.\texttt{next})$
			\State $\widetilde A = \Sort(\widehat A[0..N-|S|-1] \text{ по } x.\texttt{next})$
			
			\Comment{$i$ идёт по старым \texttt{id}, $k$ --- по $S$ в порядке \texttt{id}}
			\Comment{$j$ идёт по живым вершинам в порядке \texttt{next}, $m$ --- по $S$ в порядке \texttt{next}}
			\State $i = 0$; $j = 0$; $k = 0$; $m = 0$
			
			\While{$j < N-|S|$}
			\While{$k < |S| \ \textbf{and}\ i = S[k].\texttt{id}$}
			\State $i = i + 1$
			\State $k = k + 1$
			\Comment{пропускаем старые \texttt{id} удалённых вершин}
			\EndWhile
			
			\State $\widetilde A[j].\texttt{next} = i$
			\Comment{восстановили старый \texttt{id} вершины, на которую вёл указатель}
			
			\If{$m < |S| \ \textbf{and}\ \widetilde A[j].\texttt{next} = \widetilde S[m].\texttt{next}$}
			\Comment{значит, между предшественником и его текущим преемником надо вернуть удалённую вершину}
			\State $\widetilde A[j].\texttt{next} = \widetilde S[m].\texttt{id}$
			\State $\widetilde A[j].w = \widetilde A[j].w - \widetilde S[m].w$
			\Comment{теперь $\widetilde A[j].w$ = префиксная сумма до предшественника}
			\State $\widetilde S[m].w = \widetilde S[m].w + \widetilde A[j].w$
			\Comment{теперь $\widetilde S[m].w$ = префиксная сумма до восстановленной вершины}
			\State $m = m + 1$
			\EndIf
			
			\State $i = i + 1$
			\State $j = j + 1$
			\EndWhile
			
			\State $A = \Sort(x \in (\widetilde S \cup \widetilde A[0..N-|S|-1]) \text{ по } \texttt{id})$
			\State \Return $A$
		\end{algorithmic}
	\end{algorithm}
	
	\subsection{Оценка сложности}
	
	Обозначим через $\Sort(N)$ стоимость сортировки $N$ записей во внешней памяти.
	
	Один уровень рекурсии состоит из:
	\begin{itemize}
		\item константного числа сортировок по полям \texttt{id} и \texttt{next};
		\item константного числа линейных проходов по массивам.
	\end{itemize}
	
	Поэтому существует константа $c>0$ такая, что работа одного уровня рекурсии не превосходит
	\[
	c\,\Sort(N).
	\]
	
	\subsubsection*{Доказательство по индукции}
	
	Докажем по индукции, что
	\[
	\mathbb{E}[T(n)] \le 4c\,\Sort(n).
	\]
	
	База индукции для малых $n$ очевидна, если увеличить константу $c$ при необходимости.
	
	Пусть теперь утверждение верно для всех размеров меньше $n$.
	Тогда, раскладывая по возможным значениям $|S|$, получаем
	\[
	\mathbb{E}[T(n)]
	=
	\sum_{k=0}^{n-1}
	\Pr[|S|=k]\cdot \mathbb{E}[T(n-k)]
	+
	c\,\Sort(n).
	\]
	
	Здесь для простоты суммируем по всем $k<n$, включая $k=0$.
	
	По предположению индукции:
	\[
	\mathbb{E}[T(n)]
	\le
	\sum_{k=0}^{n-1}
	\Pr[|S|=k]\cdot 4c\,\Sort(n-k)
	+
	c\,\Sort(n).
	\]
	
	Используем стандартную оценку
	\[
	\Sort(m)=\Theta\!\left(\frac{m}{B}\log_{M/B}\frac{m}{B}\right).
	\]
	Тогда, поскольку $n-k\le n$,
	\[
	\Sort(n-k)
	\le
	\frac{n-k}{B}\log_{M/B}\frac{n}{B}
	\]
	с точностью до постоянного множителя, уже поглощённого константой $c$.
	Следовательно,
	\[
	\mathbb{E}[T(n)]
	\le
	4c\log_{M/B}\frac{n}{B}
	\sum_{k=0}^{n-1}\Pr[|S|=k]\frac{n-k}{B}
	+
	c\,\Sort(n).
	\]
	
	Но
	\[
	\sum_{k=0}^{n-1}\Pr[|S|=k](n-k)
	=
	\mathbb{E}[n-|S|].
	\]
	
	Так как каждая внутренняя вершина удаляется с вероятностью $1/4$, в среднем после сжатия остаётся не более $\frac{3n}{4}$ вершин
	(с точностью до констант, связанных с крайними вершинами списка), то есть
	\[
	\mathbb{E}[n-|S|]\le \frac{3n}{4}.
	\]
	
	Подставляя это в предыдущую оценку, получаем
	\[
	\mathbb{E}[T(n)]
	\le
	4c\log_{M/B}\frac{n}{B}\cdot \frac{3n}{4B}
	+
	c\,\Sort(n)
	\le
	3c\,\Sort(n)+c\,\Sort(n)
	=
	4c\,\Sort(n).
	\]
	
	Индукционный переход доказан, значит
	\[
	\mathbb{E}[T(n)] = O(\Sort(n)).
	\]
	
	\subsection*{Итог}
	
	Таким образом, ожидаемая сложность алгоритма равна
	\[
	\mathbb{E}[T(N)] = O(\Sort(N))
	=
	O\!\left(
	\frac{N}{B}\log_{M/B}\frac{N}{B}
	\right).
	\]
	
	\begin{Proposition}{}{}
		Пусть дано корневое дерево на $N$ вершинах.  
		У каждой вершины $v$ хранится её список детей в нужном порядке обхода:
		\[
		\texttt{struct }\{
		\ \texttt{int id;}
		\ \texttt{int kids\_num;}
		\ \texttt{int kids[kids\_num];}
		\ \}.
		\]
		
		Требуется выписать все вершины в порядке обхода DFS.
	\end{Proposition}
	
	\textbf{Идея.}
	Построим список, соответствующий эйлерову обходу дерева, и посчитаем на нём префиксные суммы весов.
	
	\subsubsection*{Построение списка}
	
	Для каждого ребра $(v,u)$, где $u$ --- ребёнок $v$, создаём две записи:
	\[
	(v\to u)
	\qquad\text{и}\qquad
	(u\to v).
	\]
	Также добавляем фиктивную стартовую запись $s$.
	
	Поле \texttt{next} задаём так, чтобы эти записи шли в порядке эйлерова обхода:
	\begin{itemize}
		\item после дуги $(v\to u)$ идём либо в первую дугу из $u$ к его первому ребёнку, либо, если детей нет, сразу в $(u\to v)$;
		\item после дуги $(u\to v)$ идём либо к следующему ребёнку вершины $v$, либо возвращаемся вверх.
	\end{itemize}
	
	Таким образом получаем список длины $2N-1$.
	
	\subsubsection*{Веса}
	
	Положим
	\[
	s.w=1,
	\]
	\[
	(v\to u).w=1 \quad \text{для дуги вниз},
	\qquad
	(u\to v).w=0 \quad \text{для дуги вверх}.
	\]
	
	Теперь запускаем алгоритм вычисления префиксных сумм на этом списке.
	
	\subsubsection*{Почему это работает}
	
	Единица добавляется ровно в момент первого входа в новую вершину.  
	Поэтому:
	\begin{itemize}
		\item корень получает номер $1$;
		\item для каждой дуги вниз $(v\to u)$ значение префиксной суммы на этой записи равно номеру вершины $u$ в порядке DFS.
	\end{itemize}
	
	Значит, если собрать пары
	\[
	(1,\text{root})
	\qquad\text{и}\qquad
	((v\to u).w,\ u)\ \text{для всех дуг вниз},
	\]
	а затем отсортировать их по первому компоненту, то вершины будут выписаны в порядке обхода DFS.
	
	\subsubsection*{Сложность}
	
	Число записей равно $O(N)$. Построение эйлерова списка и финальная сортировка требуют
	константного числа сортировок и линейных проходов, а префиксные суммы на списке считаются за $O(\Sort(N))$.
	
	Следовательно, вся задача решается за
	\[
	O(\Sort(N)).
	\]
	
	\section*{Линейное программирование}
	
	\subsection{Общая постановка задачи ЛП}
	
	\textbf{Линейное программирование (ЛП)} --- это задача оптимизации линейной функции при линейных ограничениях.
	
	В общем виде задача может быть записана так:
	\[
	\max_{x \in \mathbb{R}^n} c^T x
	\qquad \text{при} \qquad
	Ax \le b,
	\]
	где
	\[
	A \in \mathbb{R}^{m \times n}, \qquad b \in \mathbb{R}^m, \qquad c \in \mathbb{R}^n.
	\]
	
	Неравенство \(Ax \le b\) понимается \textbf{покомпонентно}, то есть
	\[
	(Ax)_i \le b_i, \qquad i=1,\dots,m.
	\]
	
	То же самое в координатной записи:
	\[
	\max \; c_1x_1+\dots+c_nx_n
	\]
	при условиях
	\[
	\begin{cases}
		a_{11}x_1+\dots+a_{1n}x_n \le b_1,\\
		a_{21}x_1+\dots+a_{2n}x_n \le b_2,\\
		\hspace{4em}\vdots\\
		a_{m1}x_1+\dots+a_{mn}x_n \le b_m.
	\end{cases}
	\]
	
 	Две наиболее употребительные записи:
	
	\[
	\max \; c^T x \quad \text{при } Ax \le b,
	\]
	--- \emph{стандартная форма}, и
	\[
	\min \; c^T x \quad \text{при } Ax = b,\; x \ge 0,
	\]
	--- \emph{каноническая форма}.
	
	Обе формы эквивалентны общей задаче после несложных преобразований.
	
	\subsection{Эквивалентные преобразования задач ЛП}
	
	Многие задачи линейного программирования можно привести к одной удобной форме.
	
	\paragraph{1) Минимизацию можно заменить максимизацией}
	\[
	\min \; c^T x
	\quad \Longleftrightarrow \quad
	\max \; (-c)^T x.
	\]
	
	\paragraph{2) Ограничение вида \texorpdfstring{$\ge$}{>=} можно заменить ограничением вида \texorpdfstring{$\le$}{<=}}
	Если
	\[
	a_k^T x \ge b_k,
	\]
	то, умножив обе части на \(-1\), получаем
	\[
	-a_k^T x \le -b_k.
	\]
	
	\paragraph{3) Равенство можно заменить двумя неравенствами}
	Если
	\[
	a_k^T x = b_k,
	\]
	то это эквивалентно системе
	\[
	\begin{cases}
		a_k^T x \le b_k,\\
		-a_k^T x \le -b_k.
	\end{cases}
	\]
	
	\paragraph{4) Свободную переменную можно заменить разностью двух неотрицательных}
	Если переменная \(x_j\) может быть любого знака, то вводят
	\[
	x_j^+ \ge 0, \qquad x_j^- \ge 0,
	\]
	и полагают
	\[
	x_j = x_j^+ - x_j^-.
	\]
	
	Это стандартный способ перейти к задаче с неотрицательными переменными.
	
	\paragraph{5) Неравенство \texorpdfstring{$\le$}{<=} можно превратить в равенство}
	Если имеется ограничение
	\[
	a_k^T x \le b_k,
	\]
	то вводится \textbf{добавочная переменная} (slack variable)
	\[
	s_k \ge 0,
	\]
	такая что
	\[
	a_k^T x + s_k = b_k.
	\]
	
	Смысл переменной \(s_k\): это \emph{остаток ресурса} или \emph{запас} до границы ограничения.
	
	\subsection{Приведение к канонической форме}
	
	Одна из самых удобных форм:
	\[
	\min \; c^T x
	\qquad \text{при} \qquad
	Ax = b, \qquad x \ge 0.
	\]
	
	Чтобы привести к ней произвольную линейную задачу, обычно делают так:
	
	\begin{enumerate}
		\item при необходимости заменяют максимум на минимум или наоборот;
		\item ограничения вида \(\ge\) заменяют на \(\le\), умножая на \(-1\);
		\item неравенства заменяют равенствами с помощью добавочных переменных;
		\item каждую свободную переменную заменяют разностью двух неотрицательных.
	\end{enumerate}
	
	\subsection{Геометрическая интерпретация}
	
	Пусть строки матрицы \(A\) равны
	\[
	a_1^T,\dots,a_m^T.
	\]
	Тогда каждое ограничение
	\[
	a_k^T x \le b_k
	\]
	задаёт \textbf{полупространство}.
	
	Следовательно, допустимое множество
	\[
	P=\{x \in \mathbb{R}^n \mid Ax \le b\}
	\]
	есть пересечение полупространств, то есть \textbf{выпуклый многогранник}.
	
	Целевая функция
	\[
	c^T x
	\]
	задаёт семейство гиперплоскостей
	\[
	c^T x = d.
	\]
	При изменении \(d\) эта гиперплоскость сдвигается параллельно самой себе. Задача ЛП состоит в том, чтобы найти такую гиперплоскость, которая:
	\begin{enumerate}
		\item ещё пересекает допустимое множество;
		\item при этом даёт наибольшее (или наименьшее) значение целевой функции.
	\end{enumerate}
	
	\begin{Example}{План производства}{}
	
	Предприятие выпускает два вида продукции:
	\[
	x_1 = \text{число столов}, \qquad x_2 = \text{число стульев}.
	\]
	
	Прибыль:
	\[
	40 \text{ у.е. за стол}, \qquad 30 \text{ у.е. за стул}.
	\]
	
	Есть два ресурса:
	\begin{itemize}
		\item древесина: не более \(18\) единиц;
		\item рабочее время: не более \(10\) часов.
	\end{itemize}
	
	На производство одной единицы продукции требуется:
	\[
	\begin{array}{c|cc}
		& \text{стол } (x_1) & \text{стул } (x_2)\\
		\hline
		\text{древесина} & 3 & 2\\
		\text{время} & 2 & 1
	\end{array}
	\]
	
	Тогда задача ЛП имеет вид:
	\[
	\max \; 40x_1 + 30x_2
	\]
	при ограничениях
	\[
	\begin{cases}
		3x_1 + 2x_2 \le 18,\\
		2x_1 + x_2 \le 10,\\
		x_1 \ge 0,\; x_2 \ge 0.
	\end{cases}
	\]
	
	\textbf{Смысл ограничений:}
	\begin{itemize}
		\item \(3x_1+2x_2 \le 18\) --- нельзя израсходовать больше древесины, чем есть;
		\item \(2x_1+x_2 \le 10\) --- нельзя превысить имеющийся запас рабочего времени;
		\item \(x_1,x_2 \ge 0\) --- отрицательное число изделий производить нельзя.
	\end{itemize}
	\end{Example}
	
	\subsection{Приближённое решение системы \texorpdfstring{$Ax \approx b$}{Ax ≈ b}}
	
	Иногда систему
	\[
	Ax=b
	\]
	нельзя решить точно: например, она переопределена или содержит ошибки измерений. Тогда ищут \(x\), для которого невязка (ошибка)
	\[
	Ax-b
	\]
	как можно меньше.
	
	Рассмотрим две постановки:
	\[
	\min_x \|Ax-b\|_2
	\qquad \text{и} \qquad
	\min_x \|Ax-b\|_1.
	\]
	
	\subsubsection*{Норма \texorpdfstring{$\ell_2$}{l2}: метод наименьших квадратов}
	
	Задача
	\[
	\min_x \|Ax-b\|_2
	\]
	имеет тот же набор решений, что и задача
	\[
	\min_x \|Ax-b\|_2^2,
	\]
	потому что квадрат --- монотонная функция на \([0,\infty)\).
	
	Именно для \textbf{квадрата нормы} выводятся \textbf{нормальные уравнения}:
	\[
	A^T A x = A^T b.
	\]
	
	\textbf{Почему получаются нормальные уравнения?}
	
	Рассмотрим функцию
	\[
	f(x)=\|Ax-b\|_2^2=(Ax-b)^T(Ax-b).
	\]
	
	Необходимое условие минимума: все частные производные по переменным
	\(x_1,\dots,x_n\) равны нулю.
	
	Вычисление этих производных приводит к равенству
	\[
	2A^T(Ax-b)=0.
	\]
	
	Следовательно,
	\[
	A^T(Ax-b)=0,
	\]
	то есть
	\[
	A^TAx=A^Tb.
	\]
	
	Это и есть нормальные уравнения.
	
	\medskip
	
	\textbf{Замечание.} Если столбцы матрицы \(A\) линейно независимы, то матрица \(A^T A\) невырождена, и решение единственно.
	
	\paragraph{Оценка сложности} для плотной матрицы \(A \in \mathbb{R}^{m\times n}\):
	\begin{itemize}
		\item вычисление \(A^T A\): \(O(mn^2)\);
		\item вычисление \(A^T b\): \(O(mn)\);
		\item решение системы размера \(n\times n\): \(O(n^3)\);
		\item память: \(O(n^2)\) для хранения \(A^T A\).
	\end{itemize}
	
	\subsubsection*{Норма \texorpdfstring{$\ell_1$}{l1}: сведение к линейному программированию}
	
	Задача
	\[
	\min_x \|Ax-b\|_1
	\]
	может быть сведена к задаче линейного программирования.
	
	Напомним, что
	\[
	\|Ax-b\|_1=\sum_{i=1}^m |(Ax-b)_i|.
	\]
	
	Чтобы убрать модуль, вводят новые переменные
	\[
	t_1,\dots,t_m \ge 0,
	\]
	где каждая \(t_i\) должна быть не меньше абсолютной величины соответствующей невязки:
	\[
	|(Ax-b)_i| \le t_i.
	\]
	
	Это эквивалентно двум линейным ограничениям:
	\[
	-(t_i) \le (Ax-b)_i \le t_i.
	\]
	
	Значит, задача переписывается так:
	\[
	\min_{x,t} \sum_{i=1}^m t_i
	\]
	при условиях
	\[
	Ax-b \le t,
	\qquad
	-(Ax-b) \le t,
	\qquad
	t \ge 0.
	\]
	
	Эквивалентно:
	\[
	\min_{x,t} \mathbf{1}^T t
	\]
	при
	\[
	\begin{cases}
		Ax - t \le b,\\
		-Ax - t \le -b,\\
		t \ge 0.
	\end{cases}
	\]
	
	\paragraph{Почему это работает}
	На оптимуме каждая переменная \(t_i\) станет минимально возможной, то есть
	\[
	t_i = |(Ax-b)_i|.
	\]
	Поэтому целевая функция
	\[
	\sum_{i=1}^m t_i
	\]
	ровно совпадает с \(\|Ax-b\|_1\).
	
	\paragraph{Размер новой задачи}
	\begin{itemize}
		\item новых переменных: \(m\);
		\item новых линейных ограничений: \(2m\);
		\item стоимость построения для плотного случая: \(O(mn)\).
	\end{itemize}
	
	\begin{Example}{ для \texorpdfstring{$\ell_1$}{l1}}{}
	Пусть
	\[
	A=
	\begin{pmatrix}
		1\\
		2
	\end{pmatrix},
	\qquad
	b=
	\begin{pmatrix}
		1\\
		4
	\end{pmatrix}.
	\]
	Тогда
	\[
	Ax-b=
	\begin{pmatrix}
		x-1\\
		2x-4
	\end{pmatrix},
	\]
	и задача
	\[
	\min_x \|Ax-b\|_1
	\]
	имеет вид
	\[
	\min_x \bigl(|x-1|+|2x-4|\bigr).
	\]
	
	Введём \(t_1,t_2 \ge 0\). Тогда получаем ЛП:
	\[
	\min_{x,t_1,t_2} t_1+t_2
	\]
	при условиях
	\[
	\begin{cases}
		x-1 \le t_1,\\
		-(x-1) \le t_1,\\
		2x-4 \le t_2,\\
		-(2x-4) \le t_2,\\
		t_1 \ge 0,\; t_2 \ge 0.
	\end{cases}
	\]
	
	Это уже обычная задача линейного программирования без модулей.
	\end{Example}
	
	\begin{Proposition}{}{}
		Как решить \texorpdfstring{$\min c^T x$}{min cTx} при \texorpdfstring{$Ax=b$}{Ax=b}, если переменные свободны по знаку	
	\end{Proposition}
	
	\subsubsection*{Идея решения}
	
	Если уже есть алгоритм, который умеет решать задачи вида
	\[
	\min \; \tilde c^T z
	\qquad \text{при} \qquad
	\tilde A z=\tilde b,\qquad z \ge 0,
	\]
	то нужно только избавиться от свободных переменных.
	
	Для каждого \(j=1,\dots,n\) вводим
	\[
	x_j=x_j^+-x_j^-,
	\qquad
	x_j^+ \ge 0,\quad x_j^- \ge 0.
	\]
	
	Тогда
	\[
	x=x^+-x^-,
	\]
	и исходная задача эквивалентна
	\[
	\min \; c^T(x^+-x^-)
	\]
	при
	\[
	A(x^+-x^-)=b,
	\qquad
	x^+ \ge 0,\quad x^- \ge 0.
	\]
	
	В матричном виде:
	\[
	\min
	\begin{pmatrix}
		c\\
		-c
	\end{pmatrix}^T
	\begin{pmatrix}
		x^+\\
		x^-
	\end{pmatrix}
	\]
	при
	\[
	\begin{pmatrix}
		A & -A
	\end{pmatrix}
	\begin{pmatrix}
		x^+\\
		x^-
	\end{pmatrix}
	=b,
	\qquad
	\begin{pmatrix}
		x^+\\
		x^-
	\end{pmatrix}
	\ge 0.
	\]
	
	Это уже задача в канонической форме.
	
	\subsubsection*{Размер новой задачи}
	\begin{itemize}
		\item было \(n\) переменных, стало \(2n\);
		\item число равенств осталось \(m\).
	\end{itemize}
	
	Построение новой задачи требует \(O(mn)\) операций для плотной матрицы.